% Generated draft: requires mathematical and source review.
\documentclass[11pt]{article}
\usepackage[margin=25mm]{geometry}
\usepackage{fontspec}
\setmainfont{Latin Modern Roman}
\usepackage{amsmath,amssymb,mathtools,unicode-math}
\setmathfont{Latin Modern Math}
\usepackage{enumitem,needspace,xurl,hyperref}
\setlength{\parindent}{0pt}
\setlength{\parskip}{5pt}
\setlength{\emergencystretch}{3em}
\newcommand{\N}{\mathbb{N}}
\newcommand{\Z}{\mathbb{Z}}
\newcommand{\Q}{\mathbb{Q}}
\newcommand{\R}{\mathbb{R}}
\newcommand{\C}{\mathbb{C}}
\newcommand{\F}{\mathbb{F}}
\newcommand{\A}{\mathbb{A}}
\newcommand{\PP}{\mathbb{P}}
\newcommand{\E}{\mathbb{E}}
\newcommand{\eps}{\varepsilon}
\newcommand{\dd}{\,\mathrm{d}}
\begin{document}
\section*{UnsolvedMath ID: 30001368}
\textbf{Source problem code:} catalogue entry
\subsection{Definitions and mathematical statement}
For a specified measurable dynamical system \(f\colon X\to X\), an invariant probability measure is a probability measure \(\mu\) satisfying \(\mu(f^{-1}(B))=\mu(B)\) for every measurable set \(B\subseteq X\). A piecewise-affine map is a map whose restriction to each member of a specified partition is affine in the relevant affine coordinates. The source does not specify the phase space, the partition, its finiteness or regularity, or the treatment of partition boundaries and discontinuities.

The source considers piecewise-affine systems preserving complementary stable and unstable cone fields. These words indicate stable and unstable tangent directions with appropriate cone invariance and hyperbolic behavior, but the record supplies no precise cone inequalities, expansion or contraction constants, invertibility assumptions, or conventions at singularities. These hypotheses therefore remain to be specified rather than inferred.

A Sinai--Ruelle--Bowen (SRB) measure is commonly described, in settings admitting the required unstable manifolds and disintegrations, as an invariant probability measure whose conditional measures on unstable manifolds are absolutely continuous with respect to the induced volume. The record does not specify the SRB convention or hypotheses ensuring that this description applies. In particular, a physical measure, defined through a basin of positive ambient volume, must not silently be substituted for an SRB measure or assumed equivalent to it.

Let \(\mathcal C\) denote the intended class of piecewise-affine systems with the stated cone-field properties, and let \(\mathcal M_{\mathrm{SRB}}(f)\) denote the SRB invariant probability measures under the intended convention. The universal question in the supplied statement is
\[
\forall f\in\mathcal C,\qquad \mathcal M_{\mathrm{SRB}}(f)\ne\varnothing.
\]
The source also asks whether SRB existence holds for ``most'' members of \(\mathcal C\), or on an open dense subset in a natural topology. Once a topology \(\tau\) on \(\mathcal C\) has been specified, the latter question can be expressed as
\[
\exists\mathcal U\subseteq\mathcal C:\quad
\mathcal U\text{ is }\tau\text{-open},\quad
\overline{\mathcal U}^{\,\tau}=\mathcal C,\quad
\forall f\in\mathcal U,\ \mathcal M_{\mathrm{SRB}}(f)\ne\varnothing.
\]
Neither \(\tau\) nor the meaning of ``most'' is supplied. Thus these are conditional formulations of the source questions, not a fully specified conjecture for a fixed class.

The supplied background literature assessment instead discusses finiteness of maximal-entropy measures for uniformly expanding or uniformly hyperbolic piecewise-affine systems. Maximal entropy and the SRB property are different requirements; that assessment does not resolve or precisely formulate the SRB questions above.
\subsection{Short English statement}
The source asks whether every piecewise-affine system preserving complementary stable and unstable cone fields has an SRB measure, and whether SRB existence holds for most systems or on an open dense set. The phase space, hyperbolicity assumptions, SRB convention, topology, and meaning of ``most'' require clarification. The supplied maximal-entropy literature assessment concerns a different measure property.
\subsection{Sources}
\begin{itemize}
\item Supplied local UnsolvedMath registry record, permanent ID \(30001368\), title \emph{SRB Measures for Piecewise-Affine Hyperbolic Systems}, source problem code \texttt{catalogue entry}, in the Oberwolfach Reports Open Problems collection. No individual proposer is identified; the supplied proposed year is 2009.
\item Supplied report citation: \emph{Mini-Workshop: Spectrum of Transfer Operators -- Recent Developments and Applications} (2010), Oberwolfach Reports, pp.\ 2699--2730, DOI \texttt{10.4171/owr/2009/49}. \url{https://doi.org/10.4171/owr/2009/49}. The literature-source label is ``original Oberwolfach report 4132.'' The supplied citation uses 2010, while the proposed year and DOI contain 2009; this discrepancy has not been resolved.
\item Supplied literature reference labeled ``Buzzi maximal-entropy theorem'': \url{https://doi.org/10.1017/S0143385708000953}. The record supplies no fuller bibliographic details or theorem hypotheses.
\item Literature assessment checked at \texttt{2026-08-21}: this is supplied dated curation, not independently verified. It labels the problem ``partially solved'' and states that finite maximal-entropy measure results are known for important classes, including piecewise-affine surface homeomorphisms, while a general statement for all uniformly expanding or uniformly hyperbolic piecewise-affine systems is not established. The background attributes this triage to supplied OWR curation and requests independent verification of time-sensitive status and cited claims. Its connection to the SRB question remains unresolved.
\item The registry labels the statement \texttt{corrected\_verified} and reports manual semantic, typographic, and LaTeX review of a meaning-preserving cleanup. These are supplied editorial qualifications, not independent verification here of the report, the mathematical formulation, or the literature status.
\end{itemize}
\end{document}
